\section{The actual-square kernel for cutoffs up to \texorpdfstring{$N^{1/20}$}{N to the 1/20}}
\label{app:kernel20}

\begingroup
\emergencystretch=3em

\paragraph{What this appendix proves and how to read it.}
This appendix proves \cref{prop:record:kernel20} for every triple
$(\beta,\kappa,\eta)$ that is admissible in the sense of
\cref{def:kernel-conditions}. It is the only estimate of the analytic stage in
which the absence of square differences in $A$ produces cancellation
(\cref{sec:kernel}). It is
applied through \cref{cor:kernel20}, at
$(\beta,\kappa,\eta)=(1/20,1/20,17/100)$, in the accepted-pair cancellation
\cref{lem:record:density-can} of \cref{sec:scale-comparison}, using part~(iii)
for the nonvanishing square classes of the finite-depth external index
and the square-clique prefix. The strict-square and fixed-prefix cases
in parts~(i) and~(ii) are special cases of this same construction. Through
the cancellation lemma the kernel also enters the wrap of the signed index
form (\cref{lem:record:wrap-expansion}).

The proof follows \cite[Prop.~6.1]{OpenAI}. As explained in
\cref{sec:kernel}, it compares a model kernel $k_0$, whose bilinear form
against $(F,J)$ is the pairing of the proposition, with a signed kernel $k$ on
the positive squares, whose bilinear form vanishes; the work is the uniform
bound $\|S-S_0\|_\infty=O(H^{-1/3})$ for their Fourier transforms
(\cref{lem:record:K20-approx}). Off the major arcs both transforms are small;
on a major arc both reduce to one arithmetic coefficient, and the two
coefficients agree exactly when the part of the denominator prime to $D$ is at
most $H$. Each error is a power of $N$ whose exponent is the left-hand side of
one of (E1)--(E4). One root-class construction
(\cref{lem:kernel-depth-classes}) supplies the arithmetic coefficient for
all three cases. We retain the square-division identity
(\cref{lem:record:nonunit-rescale}) to explain how the normalization behaves
under rescaling; the estimate itself requires no separate rescaling proof.

Besides elementary facts (geometric sums, partial summation, the Chinese
remainder theorem, Hensel lifting of square roots of units, and Plancherel's
identity on $\R/\Z$) the proof uses Dirichlet's approximation theorem: for
$Q_0\ge1$ and $\alpha\in\R$ there are coprime integers $a,q$ with
$1\le q\le Q_0$ and $|\alpha-a/q|\le1/(qQ_0)$. Weyl's inequality and the
quadratic Gauss sums at odd prime powers are proved below
(\cref{lem:record:weyl-even,lem:record:prime-power-lift}).

\paragraph{Conventions.}
The reduced denominator $q(\xi)$, the unit quadratic average $g$, the
frequency set $\Xi_{H,D}$ and the normalized transform $\widehat F$ are those
of \cref{def:kernel-notation}, and the strict square classes are those of
\cref{def:strict-square-class}; of $g$ we use its defining formula and
$|g|\le1$. Throughout, $(\beta,\kappa,\eta)$ is
admissible, and (E1)--(E4), (K0)--(K2) are the conditions of
\cref{def:kernel-conditions}. We fix $C_D,C_H\ge1$, an integer $N\ge1$,
a modulus $D\le C_DN^\kappa$ divisible by $8$, a cutoff
$C_H^{-1}N^{\beta/2}\le H\le C_HN^{\beta/2}$, and a nonvanishing square
class $r$ modulo $D$ as in \cref{prop:record:kernel20}(iii). For $p\mid D$,
$v_p(D)$ is its exponent in $D$. The root period $D'$ will be a divisor of
$D$, constructed below; all analytic bounds use only $D'\le D$.
We write $\|x\|$ for the distance
from $x\in\R$ to $\Z$, and $V_N(x)=N^{-1}\sum_{n=1}^Ne(xn)$, so that
$|V_N(x)|\le\min(1,(2N\|x\|)^{-1})$. A \emph{constant} $C$ depends only on $C_D$ and $C_H$, and ``$N$ large''
means $N\ge N_0$ for some $N_0$ depending only on $\beta$, $\kappa$, $\eta$,
$C_D$ and $C_H$; for $N$ large, $H\ge1$. As in \cref{sec:kernel}, $q$ is a
denominator and $a$, $b$, $r$ are residues or numerators (not $q=1/p$,
$r=p^{-1/2}$ or the price exponent $a$), and letters such as $A$, $B$, $J$,
$S$, $T$, $V$ have only the local meanings given here
(\cref{subsec:notation:symbols}).

\subsection{The two kernels}

The proof compares two kernels: a model kernel, built from square
root classes modulo $D$, whose pairing with $(F,J)$ is the quantity
in \cref{prop:record:kernel20}, and a signed kernel supported on
the actual positive squares, whose pairing vanishes because no
difference in question is a square. This subsection constructs the
root classes from which both kernels are built.

\begin{lemma}[Root classes at finite conductor depth]
\label{lem:kernel-depth-classes}
Let $D$ and $r$ be as above. There is a class $m_0\bmod D'$, with
$D'\mid D$, such that $\rho_D(m)=D'\one_{m\equiv m_0\,(D')}$ has mean one,
supremum $D'$ and, in the expansion $\rho_D(m)=\sum_{j\bmod D'}e(j(m-m_0)/D')$, Fourier coefficients of total absolute value $D'$.
Every $m$ in this class satisfies $m^2\equiv r\pmod D$. At a prime $p\mid D$,
the squared root variable is uniform on the residues modulo $p^k$, for every $k\ge v_p(D)$, that reduce to $r$ modulo $p^{v_p(D)}$. In particular, for $p\nmid a$, the local mean of
$e(am^2/p^k)$ is $e(ar/p^k)$ if $k\le v_p(D)$ and zero otherwise.
\end{lemma}
\begin{proof}
Fix a prime $p\mid D$. Write $h=v_p(D)$ and $v_p(r)=2j$, and choose a root $b$ of $r$
modulo $p^h$ with $v_p(b)=j$, which exists by
\cref{def:nonvanishing-square-class}. For odd $p$, restrict the root variable
to $b+p^{h-j}t$. Its square, after subtracting $r$ and dividing by
$p^h$, is
\[
 (b^2-r)/p^h+2(b/p^j)t+p^{h-2j}t^2.
\]
The linear coefficient is a unit and the quadratic coefficient is
divisible by $p$. A polynomial $u t+p v t^2+c$ with $p\nmid u$
permutes every $\Z/p^k\Z$: modulo $p$ it is bijective, and each
solution modulo $p^k$ has exactly one lift for each prescribed next
digit, since the derivative remains a unit. This proves uniformity
on the higher square residues.

At $2$, use the root class $b+2^{h-j-1}t$. The corresponding
higher-square coefficient is
\[
 (b^2-r)/2^h+(b/2^j)t+2^{h-2j-2}t^2.
\]
The linear coefficient is odd and the quadratic coefficient even,
because $h-2j\ge3$. The same digit-lifting proof applies. Replacing
$b$ by another representative of its root class does not change the
conclusion. In both cases the root class squares to $r$ at depth $h$
and is uniform over all its higher extensions; a primitive additive
character above depth $h$ therefore has mean zero.

The root periods are $p^{h-j}$ at odd primes and $2^{h-j-1}$ at
$2$. Their product $D'$ divides $D$. The Chinese remainder theorem
gives one root class $b_0\bmod D'$, with normalized weight
$D'\one_{b=b_0\,(D')}$. This weight has mean one, supremum $D'$,
and exactly $D'$ Fourier coefficients of modulus one. The local character
means follow by averaging over the uniform higher square residues: for $k>h$ the factor is
$e(ar/p^k)\E_{t\bmod p^{k-h}}e(at/p^{k-h})=0$.
\end{proof}

\begin{definition}[Root weight, unit sieve and the two kernels]
\label{def:record:K20-kernels}
Use the root weight $\rho_D$ and period $D'$ of
\cref{lem:kernel-depth-classes}. For an odd prime
$p\nmid D$ put
\[
 w_p(m)=\frac{p\one_{p\mid m}-1}{p-1}=\frac1{p-1}\sum_{j=1}^{p-1}e(jm/p),
\]
and $w_u=\prod_{p\mid u}w_p$ for squarefree $u$ coprime to $D$ ($w_1=1$).
With $u$ running over the squarefree integers $1\le u\le H$ coprime to $D$,
define
\begin{align*}
 k_0(n)&=\frac DN\one_{1\le n\le N}\one_{n\equiv r\,(D)}
  \sum_{\xi\in\Xi_{H,D}}g(\xi)e(-\xi n),\\
 k(n)&=\sum_{1\le m\le\sqrt N}\one_{n=m^2}\frac{2m}N\,\rho_D(m)
  \sum_u(-1)^{\omega(u)}w_u(m),
\end{align*}
and $S_0(\alpha)=\sum_nk_0(n)e(\alpha n)$, $S(\alpha)=\sum_nk(n)e(\alpha n)$.
These objects have the following properties.
\begin{enumerate}[label=(\alph*)]
\item $\rho_D$ has mean one over its period $D'$, $\|\rho_D\|_\infty=D'$,
 $\rho_D(m)=\sum_{j\bmod D'}e(j(m-m_0)/D')$, and $m^2\equiv r\pmod D$
 whenever $\rho_D(m)\ne0$.
\item $|w_u|\le1$; $w_p$ has mean zero modulo $p$; the Fourier coefficients
 of $w_u$ have total absolute value one; and for $q$ coprime to $D$
 \begin{equation}\label{eq:record:unitsieve}
  \sum_{u\mid\rad(q)}(-1)^{\omega(u)}w_u(m)=\prod_{p\mid q}\big(1-w_p(m)\big)
  =\frac q{\varphi(q)}\one_{(m,q)=1}.
 \end{equation}
\item $S(\alpha)=\sum_{m\le\sqrt N}\frac{2m}N\Phi(m)e(\alpha m^2)$, where
 $\Phi(m)=\rho_D(m)\sum_u(-1)^{\omega(u)}w_u(m)$ is a finite sum of additive
 characters in $m$ with coefficients of total absolute value at most $D'H$.
\end{enumerate}
\end{definition}

\begin{proof}[Proof of (a)--(c)]
(a) The root-class lemma gives the mean, supremum and square condition;
the Fourier formula is orthogonality modulo $D'$.
(b) $w_p(m)$ is $1$ if $p\mid m$ and $-1/(p-1)$ otherwise, and
$1-w_p(m)=\frac p{p-1}\one_{p\nmid m}$; expanding the product gives
\eqref{eq:record:unitsieve}, since $\prod_{p\mid q}p/(p-1)=q/\varphi(q)$.
(c) The total absolute value of the coefficients is submultiplicative under
products; $\rho_D$ contributes $D'$, each $w_u$ one, and there are at most $H$
values of $u$.
\end{proof}

\paragraph{How the kernels are matched.}
By (a), $k$ lives on the squares in the class $r$ modulo $D$, the class of
$k_0$, and the factor $2m/N$ is the discrete form of $dn=2m\,dm$, matching the
weight $1/N$ of $k_0$. The divisor sum \eqref{eq:record:unitsieve} is the
density of the uniform law on the units modulo $q$; it turns averages of
$e(\xi m^2)$ over residues into averages over units, which is how $g$ arises
(\cref{lem:record:kernel-coefficient}). It is truncated at $u\le H$ so that
$\Phi$ has Fourier coefficients of total size at most $D'H$. The same
construction in the strict-square case underlies \cite[Prop.~6.1]{OpenAI}.

\begin{lemma}[The model kernel as a sum over centres]
\label{lem:record:K20-centres}
We have
\[
 S_0(\alpha)=\sum_{\xi\in\Xi_{H,D}}\sum_{j\bmod D}g(\xi)e(-jr/D)
 V_N(\alpha-\xi+j/D).
\]
The \emph{centres} $\xi-j/D$ are distinct, there are at most $DH^2$ of them,
each is a reduced fraction of denominator at most $DH$, and each coefficient
has modulus at most one. Every $\alpha_0\in\Q/\Z$ is uniquely
$\alpha_0=\alpha_D+\alpha_1$, where the prime factors of $q(\alpha_D)$ divide
$D$ and $q_1=q(\alpha_1)$ is coprime to $D$; the point $\alpha_0$ is a centre
exactly when $q(\alpha_D)\mid D$ and $\alpha_1\in\Xi_{H,D}$, and its
coefficient is then $c_0(\alpha_0)=g(\alpha_1)e(\alpha_Dr)$. Put
$c_0(\alpha_0)=0$ otherwise.
\end{lemma}

\begin{proof}
Expand $\one_{n\equiv r\,(D)}=D^{-1}\sum_{j\bmod D}e(j(n-r)/D)$. If
$\xi-j/D=\xi'-j'/D$, then $\xi-\xi'$ has denominator dividing
$q(\xi)q(\xi')$, which is coprime to $D$, and also dividing $D$; so
$\xi=\xi'$ and $j\equiv j'$. The number of centres is at most
$D\sum_{q\le H}\varphi(q)\le DH^2$. The decomposition of $\alpha_0$ is the
partial fraction decomposition; a centre $\xi-j/D$ has $\alpha_1=\xi$ and
$\alpha_D=-j/D$, whence the coefficient $g(\xi)e(-jr/D)=g(\alpha_1)e(\alpha_Dr)$.
\end{proof}

\subsection{Weyl's inequality with even denominators}

The minor-arc estimate needs a bound for quadratic exponential
sums with explicit constants and with no parity condition on the
denominator, because even denominators occur through the modulus
$D$. We prove the form needed.

\begin{lemma}[Weyl's inequality with a linear twist]
\label{lem:record:weyl-even}
Let $\alpha\in\R$ and let $a/q$ be a reduced fraction with $q\ge1$ and
$|\alpha-a/q|\le q^{-2}$. For all $\lambda\in\R$ and all integers
$1\le M'\le M$,
\[
 \Big|\sum_{m=1}^{M'}e(\alpha m^2+\lambda m)\Big|^2
 \le M+200\Big(\frac Mq+1\Big)\big(M+q\log(2q)\big).
\]
No parity condition on $q$ is needed.
\end{lemma}

This is Weyl's inequality with explicit constants (compare
\cite[Lemma~2.4]{Vaughan97}); in the proof, two indices $h,h'$ in one block of $\Delta=\max(1,\lfloor q/4\rfloor)$ consecutive integers differ by $j=|h-h'|$ with $0<2j<q/2$ when $q\ge4$, so $q\nmid2aj$ whatever the parity of $q$.

\begin{proof}
Squaring and grouping the pairs $(m,m+h)$ by $h$,
\[
 \Big|\sum_{m\le M'}e(\alpha m^2+\lambda m)\Big|^2\le M'+2\sum_{h=1}^{M'-1}
 \Big|\sum_{m\in I_h}e(2\alpha hm)\Big|,
\]
where $I_h$ is an interval of length at most $M'$; the phases
$e(\alpha h^2+\lambda h)$ are unimodular, so the twist $\lambda$ plays no
role. A geometric sum satisfies $|\sum_{m\in I}e(xm)|\le\min(|I|,(2\|x\|)^{-1})$,
so the left-hand side is at most $M+2\sum_{h=1}^M\min(M,\|2\alpha h\|^{-1})$.

Let $\Delta=\max(1,\lfloor q/4\rfloor)$ and split $\{1,\dots,M\}$ into at most
$M/\Delta+1$ blocks of $\Delta$ consecutive integers. Suppose $q\ge4$, and let
$h\ne h'$ lie in one block, so $j=|h-h'|$ satisfies $0<j<\Delta\le q/4$. Then
$q\nmid2j$, hence $q\nmid2aj$, so $\|2aj/q\|\ge1/q$; and
$|2j(\alpha-a/q)|\le2j/q^2<1/(2q)$. Therefore
$\|2\alpha h-2\alpha h'\|=\|2\alpha j\|>1/(2q)$: the points $2\alpha h$ of one
block are $1/(2q)$-separated on $\R/\Z$. Each set
$\{x:k/(2q)\le\|x\|<(k+1)/(2q)\}$ is the union of two half-open intervals of
length $1/(2q)$ and contains at most two points of a block. The points with
$k=0$ contribute at most $2M$, the others at most $2\cdot2q/k$ for each $k$,
so a block contributes at most $2M+4q\sum_{k=1}^q1/k\le2M+12q\log(2q)$. Since
$\Delta\ge q/8$ for $q\ge4$, there are at most $8M/q+1$ blocks, and
$2\sum_h\min(M,\|2\alpha h\|^{-1})\le2(8M/q+1)(2M+12q\log(2q))
\le192(M/q+1)(M+q\log(2q))$. If $q\le3$, then $2\sum_h\min(M,\|2\alpha h\|^{-1})\le2M^2\le6(M/q)M$, which is
smaller.
\end{proof}

\subsection{Arcs}

For a reduced fraction $a/q$ with $1\le q\le N^\eta$ put
$\mathfrak M(a/q)=\{\alpha\in\R/\Z:\|\alpha-a/q\|\le2N^{-1+\eta}\}$, and let
$\mathfrak m$ be the complement of the union of these major arcs. For
$\alpha_0=a/q$ and every squarefree $1\le u\le H$ coprime to $D$ let
$a_u(m)=\rho_D(m)w_u(m)e(\alpha_0m^2)$, let $A_u$ be its mean over any one period (the value is independent of the period chosen),
and put
\[
 c(\alpha_0)=\sum_u(-1)^{\omega(u)}A_u .
\]

\begin{lemma}[Arcs]\label{lem:record:arcs}
For $N$ large the following hold.
\begin{enumerate}[label=(\roman*)]
\item The major arcs are pairwise disjoint. Every $\alpha\in\mathfrak m$ has a
 reduced fraction $a/q$ with $N^\eta<q\le N^{1-\eta}$ and $|\alpha-a/q|\le q^{-2}$.
\item Every centre of $S_0$ is the centre of a major arc. If $\alpha=\alpha_0+y$
 with $\alpha_0=a/q$, $q\le N^\eta$ and $|y|\le2N^{-1+\eta}$, then every centre
 $\alpha_c\ne\alpha_0$ satisfies $\|\alpha-\alpha_c\|\ge(2qDH)^{-1}$.
\item For $\alpha\in\mathfrak m$,
 $|S(\alpha)|\le CN^{\kappa+\beta/2-\eta/2}(\log2N)^{1/2}$ and
 $|S_0(\alpha)|\le CN^{\kappa+\beta-\eta}$.
\item For $\alpha$ as in (ii), $S_0(\alpha)=c_0(\alpha_0)V_N(y)+R_0$ with
 $|R_0|\le CN^{\eta+2\kappa+3\beta/2-1}$, and $S(\alpha)=c(\alpha_0)V_N(y)+R_1$
 with $|R_1|\le C(N^{2\eta+2\kappa+\beta-1/2}+N^{\eta+\beta/2-1})$.
\end{enumerate}
The exponents $\kappa+\beta/2-\eta/2$, $\kappa+\beta-\eta$,
$\eta+2\kappa+3\beta/2-1$, $2\eta+2\kappa+\beta-1/2$ are the left-hand sides of
(E1)--(E4), hence below $-\beta/6$, and $\eta+\beta/2-1$ is smaller than the
third. At the triple $(1/20,1/20,17/100)$ of \cref{cor:kernel20} the five
exponents are $-1/100$, $-7/100$, $-131/200$, $-1/100$, $-161/200$, below
$-1/120$, and the seven kernel slacks of (E1)--(E4), (K0)--(K2), computed in
\cref{cor:kernel20}, are $1/600$, $37/600$, $97/150$, $1/600$, $49/100$,
$19/200$, $117/200$.
\end{lemma}

\begin{proof}
We use $D'\le D\le C_DN^\kappa$ and $H\le C_HN^{\beta/2}$, so
$DH\le C_DC_HN^{\kappa+\beta/2}$.

(i) Two distinct reduced fractions of denominators at most $N^\eta$ are at
distance at least $N^{-2\eta}$, and $4N^{-1+\eta}<N^{-2\eta}$ for $N$ large, by
(K0). For $\alpha\in\mathfrak m$, Dirichlet's theorem with
$Q_0=\lfloor N^{1-\eta}\rfloor$ gives a reduced $a/q$ with $q\le Q_0$ and
$|\alpha-a/q|\le1/(qQ_0)\le q^{-2}$. If $q\le N^\eta$, then
$|\alpha-a/q|\le1/Q_0\le2N^{-1+\eta}$, so $\alpha\in\mathfrak M(a/q)$, which is
impossible. Hence $N^\eta<q\le N^{1-\eta}$.

(ii) Every centre has denominator at most
$DH\le C_DC_HN^{\kappa+\beta/2}<N^\eta$ for $N$ large, by (K1). Two distinct
fractions of denominators $q$ and at most $DH$ are at distance at least
$(qDH)^{-1}$. Since $4qDHN^{-1+\eta}\le4C_DC_HN^{2\eta+\kappa+\beta/2-1}\le1$
for $N$ large, by (K2), we have $|y|\le(2qDH)^{-1}$, and the triangle
inequality gives the claim.

(iii) Let $\alpha\in\mathfrak m$, with $a/q$ from (i), and
$M=\lfloor\sqrt N\rfloor$. By (K0), $\eta<1/2$, so $M\le N^{1/2}\le N^{1-\eta}$;
also $M^2/q\le N^{1-\eta}$ and $q\le N^{1-\eta}$. By
\cref{lem:record:weyl-even}, uniformly in $\lambda$ and $M'\le M$, expanding
the product,
\[
 \Big|\sum_{m\le M'}e(\alpha m^2+\lambda m)\Big|^2\le M+200\Big(\frac{M^2}q
 +M\log2q+M+q\log2q\Big)\le801N^{1-\eta}\log(2N).
\]
The weights $2m/N$ increase and are at most $2/\sqrt N$, so partial summation
bounds $|\sum_{m\le\sqrt N}\frac{2m}Ne(\alpha m^2+\lambda m)|$ by $4/\sqrt N$
times the maximum over $M'\le M$ of the sums above, that is, by
$116N^{-\eta/2}(\log2N)^{1/2}$. Expanding $\Phi$ by
\cref{def:record:K20-kernels}(c),
$|S(\alpha)|\le116D'HN^{-\eta/2}(\log2N)^{1/2}
\le116C_DC_HN^{\kappa+\beta/2-\eta/2}(\log2N)^{1/2}$.
By (ii), every centre is the centre of a major arc, so it lies at distance
more than $2N^{-1+\eta}$ from $\alpha$; each term of
\cref{lem:record:K20-centres} is therefore at most $(4N^\eta)^{-1}$ in modulus,
and $|S_0(\alpha)|\le DH^2/(4N^\eta)\le C_DC_H^2N^{\kappa+\beta-\eta}$.

(iv) \emph{The model kernel.} If $\alpha_0$ is a centre, its term in
\cref{lem:record:K20-centres} is $c_0(\alpha_0)V_N(y)$; if not,
$c_0(\alpha_0)=0$. By (ii), each of the at most $DH^2$ other centres has
$|V_N(\alpha-\alpha_c)|\le qDH/N$, so
$|R_0|\le qD^2H^3/N\le C_D^2C_H^3N^{\eta+2\kappa+3\beta/2-1}$.

\emph{The actual kernel.} Since $(m+q)^2-m^2=q(2m+q)$, the factor
$e(\alpha_0m^2)$ has period $q$, so $a_u$ has a period $T_u$ dividing
$\lcm(D',u,q)$, and $T_u\le D'uq$. Also $|a_u|\le D'$ and
$|A_u|\le\E|\rho_D|=1$. Hence $E_u(x)=\sum_{1\le m\le x}a_u(m)-A_ux$
satisfies $|E_u(x)|\le D'T_u+T_u\le2D'T_u$ for real $x\ge0$: the complete
periods contribute exactly $A_u$ times their length, and fewer than $T_u$
terms remain. Put $X=\sqrt N$ and $f_y(x)=\frac{2x}Ne(yx^2)$. Since
$e(\alpha m^2)=e(\alpha_0m^2)e(ym^2)$, \cref{def:record:K20-kernels}(c) gives
$S(\alpha)=\sum_u(-1)^{\omega(u)}\sum_{m\le X}f_y(m)a_u(m)$. Stieltjes
integration by parts against $\sum_{m\le x}a_u(m)=A_ux+E_u(x)$, with
$f_y(0)=E_u(0)=0$, gives
\[
 \sum_{m\le X}f_y(m)a_u(m)=A_u\int_0^Xf_y(x)\,dx+f_y(X)E_u(X)
 -\int_0^XE_u(x)f_y'(x)\,dx .
\]
Here $|f_y(X)|=2/\sqrt N$ and
$\int_0^X|f_y'|\le\int_0^X(2/N+8\pi|y|x^2/N)\,dx\le2/\sqrt N+9|y|\sqrt N$, so
the last two terms are at most $18D'T_u(1+|y|N)/\sqrt N$. The substitution
$v=x^2$ gives $\int_0^Xf_y=N^{-1}\int_0^Ne(yv)\,dv$, which differs from
$V_N(y)$ by at most $2\pi|y|$, because $|e(yn)-e(yv)|\le2\pi|y|$ for
$v\in[n-1,n]$. Summing over the at most $H$ values of $u$, with
$\sum_{u\le H}u\le H^2$, $T_u\le D'uq$ and $|A_u|\le1$,
\[
 S(\alpha)=c(\alpha_0)V_N(y)+R_1,\qquad
 |R_1|\le\frac{18qD^2H^2(1+|y|N)}{\sqrt N}+2\pi|y|H .
\]
With $q\le N^\eta$, $1+|y|N\le3N^\eta$ and $|y|\le2N^{-1+\eta}$ this is at most
\[
 54C_D^2C_H^2N^{2\eta+2\kappa+\beta-1/2}+4\pi C_HN^{\eta+\beta/2-1}.\qedhere
\]
\end{proof}

\subsection{Local factors and the arithmetic coefficient}

For $p\mid D$, let $\mathcal C_p$ be the residue class $b+p^{h-j}\Z$ (resp.\ $b+2^{h-j-1}\Z$ at $p=2$) constructed in the proof of \cref{lem:kernel-depth-classes}, with $h,j,b$ as there. The weight $\rho_D$ is the product of their
normalized indicators. We collect the two local rules needed for the
arithmetic coefficient: the root-class rule at primes of $D$, and the
quadratic Gauss sums at the remaining primes.

\begin{lemma}[Local means at the primes of $D$, and Gauss sums]
\label{lem:record:prime-power-lift}
\begin{enumerate}[label=(\alph*)]
\item Let $p\mid D$, $j\ge1$ and $p\nmid a$. The mean of $e(am^2/p^j)$ over the
 residues $m\in\mathcal C_p$ modulo $p^{\max(j,v_p(D))}$ is $e(ar/p^j)$ if
 $j\le v_p(D)$, and $0$ if $j>v_p(D)$.
\item Let $p\nmid D$ be an odd prime, $j\ge1$ and $p\nmid a$, and put
 $G_{p^j}(a)=\sum_{m\bmod p^j}e(am^2/p^j)$. Then $|G_{p^j}(a)|=p^{j/2}$. If
 $j\ge2$, then $\sum_{m\in(\Z/p^j\Z)^\times}e(am^2/p^j)=0$. The twisted sum
 \[
  \Gamma_j=p^{-j}\sum_{m\bmod p^j}w_p(m)e(am^2/p^j)
  \quad\text{satisfies}\quad|\Gamma_j|\le2p^{-j/2}.
 \]
\end{enumerate}
\end{lemma}

\begin{proof}
(a) This is the local character mean of \cref{lem:kernel-depth-classes}.

(b) Writing the pairs $(m,m')$ as $(m,m+h)$,
\[
 |G_{p^j}(a)|^2=\sum_{h\bmod p^j}e(ah^2/p^j)\sum_{m\bmod p^j}e(2ahm/p^j)=p^j,
\]
because the inner sum is $p^j$ if $p^j\mid h$ and $0$ otherwise, as
$p\nmid2a$. Let $j\ge2$, and write $m=x+p^{j-1}h$ with $0\le x<p^{j-1}$ and
$0\le h<p$; then $m^2\equiv x^2+2xhp^{j-1}\pmod{p^j}$, since $2j-2\ge j$.
Summing over $h$ removes the $x$ with $p\nmid x$, while for $p\mid x$ the $p$
lifts $m$ of $x$ are divisible by $p$ and have $m^2\equiv x^2\pmod{p^j}$. Hence
$G_{p^j}(a)$ and the sum of $e(am^2/p^j)$ over the $m$ divisible by $p$ both
equal $p\sum_{0\le x<p^{j-1},\,p\mid x}e(ax^2/p^j)$, and the sum over the
units, which is their difference, vanishes. For $\Gamma_j$ write
$w_p(m)=\frac p{p-1}\one_{p\mid m}-\frac1{p-1}$. If $j=1$, the first part
contributes the single term $m=0$, of size $\frac1{p-1}$, and the second has
modulus $\frac1{p-1}p^{-1/2}$; so $|\Gamma_1|\le\frac{1+p^{-1/2}}{p-1}\le2p^{-1/2}$
for $p\ge3$. If $j\ge2$, the sum over the $m$ divisible by $p$ equals
$G_{p^j}(a)$, so $\Gamma_j=p^{-j}G_{p^j}(a)$ and $|\Gamma_j|=p^{-j/2}$.
\end{proof}

\begin{lemma}[The arithmetic coefficient]\label{lem:record:kernel-coefficient}
Let $\alpha_0=a/q$ be reduced and $\alpha_0=\alpha_D+\alpha_1$ as in
\cref{lem:record:K20-centres}, with $q_1=q(\alpha_1)$. If $q_1\le H$, then
$c(\alpha_0)=c_0(\alpha_0)$. If $q_1>H$, then $c_0(\alpha_0)=0$ and
\[
 |c(\alpha_0)|\le3^{\omega(q_1)}q_1^{-1/2}\le3^{\pi(729)}q_1^{-1/3}
 <3^{\pi(729)}H^{-1/3}.
\]
\end{lemma}

\begin{proof}
Let $m$ be uniform modulo $\lcm(D',u,q)$, a common period of the factors of
$a_u$. By the Chinese remainder theorem, its residues modulo the part of this
modulus at the primes of $D$ and modulo the part prime to $D$ are independent
and uniform. Now $\rho_D(m)$ and $e(\alpha_Dm^2)$ depend only on the first,
and $w_u(m)$ and $e(\alpha_1m^2)$ only on the second. Hence
$A_u=A^{(D)}A^{(1)}_u$ with $A^{(D)}=\E_m\rho_D(m)e(\alpha_Dm^2)$ and
$A^{(1)}_u=\E_mw_u(m)e(\alpha_1m^2)$, and both factor over prime powers.

\emph{The $D$-part.} Write $\alpha_D=\sum_{p\mid D}a_p/p^{j_p}$ with
$p\nmid a_p$ when $j_p\ge1$. Since $\rho_D$ is the product of the normalized
indicators of the root classes, the local factor at a prime $p\mid D$ with
$j_p\ge1$ is the mean of $e(a_pm^2/p^{j_p})$ over $\mathcal C_p$, which by
\cref{lem:record:prime-power-lift}(a) is $e(a_pr/p^{j_p})$ if
$j_p\le v_p(D)$ and $0$ otherwise; a prime of $D$ with $j_p=0$ contributes
$1$. Therefore $A^{(D)}=e(\alpha_Dr)$ if $q(\alpha_D)\mid D$, and $A^{(D)}=0$
otherwise.

\emph{The coprime part.} The denominator $q_1$ is odd, being prime to $8$. If
$p\mid u$ and $p\nmid q_1$, the local factor at $p$ is the mean of $w_p$,
which is zero. So $A_u=0$ unless $u\mid\rad(q_1)$.

\emph{Case $q_1\le H$.} Every divisor of $\rad(q_1)$ is squarefree, coprime
to $D$ and at most $H$, so all of them occur among the $u$, and by
\eqref{eq:record:unitsieve}
\[
 c(\alpha_0)=A^{(D)}\,\E_{m\bmod q_1}\Big[\frac{q_1}{\varphi(q_1)}
 \one_{(m,q_1)=1}e(\alpha_1m^2)\Big]
 =\frac{A^{(D)}}{\varphi(q_1)}\sum_{m\in(\Z/q_1\Z)^\times}e(\alpha_1m^2).
\]
If $q_1$ is squarefree, the last average is $g(\alpha_1)$
(\cref{def:kernel-notation}), and $c(\alpha_0)=c_0(\alpha_0)$ by
\cref{lem:record:K20-centres}; both vanish unless $q(\alpha_D)\mid D$. If $q_1$
is not squarefree, the unit sum factors over the prime powers $p^j\,\|\,q_1$
and vanishes at one with $j\ge2$ (\cref{lem:record:prime-power-lift}(b)); then
$c(\alpha_0)=0=c_0(\alpha_0)$.

\emph{Case $q_1>H$.} Then $\alpha_1\notin\Xi_{H,D}$ and $c_0(\alpha_0)=0$. Let
$u\mid\rad(q_1)$ and $p^j\,\|\,q_1$, and let $a'/p^j$ be the $p$-component of
$\alpha_1$. The local factor of $A^{(1)}_u$ at $p$ is the normalized Gauss sum
$p^{-j}G_{p^j}(a')$, of modulus $p^{-j/2}$, if $p\nmid u$, and $\Gamma_j$, of
modulus at most $2p^{-j/2}$, if $p\mid u$ (\cref{lem:record:prime-power-lift}(b)).
So $|A^{(1)}_u|\le2^{\omega(u)}q_1^{-1/2}$, and with $|A^{(D)}|\le1$,
\[
 |c(\alpha_0)|\le q_1^{-1/2}\sum_{u\mid\rad(q_1)}2^{\omega(u)}=3^{\omega(q_1)}q_1^{-1/2}.
\]
Every prime $p\ge729=3^6$ has $p^{1/6}\ge3$, and at most $\pi(729)$ primes are
smaller; so $3^{\omega(q)}\le3^{\pi(729)}q^{1/6}$ for every $q\ge1$.
\end{proof}

\subsection{The uniform approximation and the kernel estimate}

All the arc estimates are now available. The lemma combines them
into the uniform bound $\|S-S_0\|_\infty=O(H^{-1/3})$ for the
Fourier transforms of the two kernels, and
\cref{prop:record:kernel20} then follows by Plancherel.

\begin{lemma}[Uniform Fourier approximation]\label{lem:record:K20-approx}
For $N$ large,
$\sup_{\alpha\in\R/\Z}|S(\alpha)-S_0(\alpha)|\le(3^{\pi(729)}+4)H^{-1/3}$.
\end{lemma}

\begin{proof}
Let $s_{\min}>0$ be the least of the slacks of (E1)--(E4). By
\cref{lem:record:arcs}, every power of $N$ in its parts (iii) and (iv) has
exponent at most $-\beta/6-s_{\min}$, while $H^{-1/3}\ge C_H^{-1/3}N^{-\beta/6}$.
On $\mathfrak m$, \cref{lem:record:arcs}(iii) bounds $|S|+|S_0|$. On a major
arc, \cref{lem:record:arcs}(iv), $|V_N(y)|\le1$ and
\cref{lem:record:kernel-coefficient} give
$|S-S_0|\le3^{\pi(729)}H^{-1/3}+|R_0|+|R_1|$. Each of the four quantities $|S|,|S_0|$ in \cref{lem:record:arcs}(iii) and $|R_0|,|R_1|$ in \cref{lem:record:arcs}(iv) is therefore at most $H^{-1/3}$ for $N$ large, the factor
$(\log2N)^{1/2}$ included.
\end{proof}

\begin{proof}[Proof of \cref{prop:record:kernel20}]
First prove (iii), using \cref{def:record:K20-kernels} with the nonvanishing
square class $r\equiv b-a\pmod D$. Extend $F,J$ by zero, and for a finitely
supported $h$ put
$B_h(F,J)=N^{-1}\sum_{x,n}h(n)F(x)J(x+n)$. With
$\widetilde F(\alpha)=\sum_xF(x)e(-x\alpha)$ and
$S_h(\alpha)=\sum_nh(n)e(\alpha n)$, so that $S_{k_0}=S_0$ and $S_k=S$,
orthogonality and then Cauchy--Schwarz and Plancherel give
\[
 B_h(F,J)=\frac1N\int_0^1S_h(\alpha)\widetilde F(-\alpha)\widetilde J(\alpha)\,d\alpha,
 \qquad
 |B_{k-k_0}(F,J)|\le\frac{\|S-S_0\|_\infty\|F\|_{\ell^2}\|J\|_{\ell^2}}N.
\]
Every term of $B_k(F,J)$ contains $F(x)J(x+m^2)$ with $m\ge1$, which vanishes
because no difference of the supports is a positive square; so
$B_k(F,J)=0$. This is the step at which square-difference-freeness is used
(\cref{sec:kernel}). If $F(x)J(y)\ne0$, then $1\le y-x\le N-1$ by the hypothesis of \cref{prop:record:kernel20} that every point of $\supp F$ is smaller than every point of $\supp J$,
and $y-x\equiv b-a\equiv r\pmod D$, so both indicators in $k_0(y-x)$ equal one
and
\[
 B_{k_0}(F,J)=\frac D{N^2}\sum_{x,y}F(x)J(y)\sum_{\xi\in\Xi_{H,D}}g(\xi)
 e(-\xi(y-x))=D\sum_{\xi\in\Xi_{H,D}}g(\xi)\widehat F(-\xi)\widehat J(\xi).
\]
Hence the left-hand side of \eqref{eq:record:K20} is
$|B_{k_0-k}(F,J)|\le(3^{\pi(729)}+4)H^{-1/3}\|F\|_{\ell^2}\|J\|_{\ell^2}/N$ by
\cref{lem:record:K20-approx}, and $H^{-1/3}\le C_H^{1/3}N^{-\beta/6}$; so
\eqref{eq:record:K20} holds with $C=(3^{\pi(729)}+4)C_H^{1/3}$ for $N$ large.
Part (i) is the case of unit square classes and $v_2(D)=3$; the valuation
conditions of (ii) also describe nonvanishing square classes, so the same
bound applies there. Its rescaling identity is proved next. Finally, a
function bounded by one on one class modulo $D$ in an interval $I$ has
$\|F\|_{\ell^2}^2\le|I|/D+1$, giving the last assertion.
\end{proof}

\subsection{Nonunit differences}

This subsection is explanatory and is not used in the estimate
just proved. It records exactly how the kernel statement and its
normalization transform when both supports are divided by a common
square $R$, which is the rescaling used in the main text when the
forced difference is not a unit modulo $D$.

\begin{lemma}[Division by an even-valuation square]
\label{lem:record:nonunit-rescale}
In the setting of \cref{prop:record:kernel20}(ii), let $\mathcal P_0'$ be the
set of $p\in\mathcal P_0$ with $v_p(b-a)=2j_p\ge2$, and put
$R=\prod_{p\in\mathcal P_0'}p^{2j_p}$, $D_R=D/R$ and $N'=\lceil N/R\rceil$.
Then $a\equiv b\pmod R$. Let $a_R\in(-R,0]$ with $a_R\equiv a\pmod R$, and
$F'(n')=F(a_R+Rn')$, $J'(n')=J(a_R+Rn')$ for $n'\in\Z$.
\begin{enumerate}[label=(\alph*)]
\item $F'$ and $J'$ are supported in $[1,N']$, in the classes
 $a'=(a-a_R)/R$ and $b'=(b-a_R)/R$ modulo $D_R$, and
 $\|F'\|_{\ell^2}=\|F\|_{\ell^2}$, $\|J'\|_{\ell^2}=\|J\|_{\ell^2}$. Moreover
 $D_R=8d_R$ with $d_R=d/R$ odd, $D_R$ and $D$ have the same prime divisors,
 and $b'-a'$ is a strict square class modulo $D_R$.
\item Every point of $\supp F'$ is smaller than every point of $\supp J'$, and
 no such difference is the square of a positive integer.
\item $\Xi_{H,D_R}=\Xi_{H,D}$, and $\xi\mapsto R\xi$ is a bijection of this set
 onto itself that preserves conductors, with $g(R\xi)=g(\xi)$.
\item With $\widehat F_N$ normalized by $N$ and $\widehat{F'}_{N'}$ by $N'$,
 \[
  D\sum_{\xi\in\Xi_{H,D}}g(\xi)\widehat F_N(-\xi)\widehat J_N(\xi)
  =R\Big(\frac{N'}N\Big)^2\,D_R\sum_{\xi'\in\Xi_{H,D_R}}g(\xi')
  \widehat{F'}_{N'}(-\xi')\widehat{J'}_{N'}(\xi').
 \]
\end{enumerate}
\end{lemma}

\begin{proof}
Since $p^{2j_p}$ divides $b-a$ for $p\in\mathcal P_0'$, $R\mid b-a$; and
$R\mid D$, because $2j_p<2m_p$.

(a) If $F(n)\ne0$, then $1\le n\le N$ and $n\equiv a\equiv a_R\pmod R$, so
$n=a_R+Rn'$ with $n'=(n-a_R)/R$. As $1-a_R\in[1,R]$, $n'\ge1/R>0$; and
$n'\le(N+R-1)/R$, so $n'\le\lceil N/R\rceil$. The congruence $n\equiv a\pmod D$
is equivalent to $Rn'\equiv a-a_R\pmod D$, that is, to $n'\equiv a'\pmod{D_R}$,
because $R$ divides both $a-a_R$ and $D$. The same holds for $J$, and the map
$n\mapsto n'$ is a bijection of supports, which preserves the $\ell^2$ norms.
$R$ is an odd square, so $d_R$ is odd and $R\equiv1\pmod8$. At
$p\in\mathcal P_0'$ the exponent of $p$ in $D_R$ is $2m_p-2j_p\ge2$; at every
other prime the exponent is unchanged. Now $b'-a'\equiv(b-a)/R\pmod{D_R}$. It
is $\equiv1\pmod8$, since $b-a\equiv1$ and $R\equiv1$. At $p\in\mathcal P_0'$,
$(b-a)/R$ is the unit part $(b-a)/p^{2j_p}$, a nonzero square modulo $p$,
divided by $R/p^{2j_p}$, the square of an integer prime to $p$; so it is a
nonzero square modulo $p$. At every other odd prime of $D_R$, $b-a$ is a
nonzero square modulo $p$ and $R$ is the square of an integer prime to $p$.
So $b'-a'$ is a strict square class modulo $D_R$.

(b) The map $n\mapsto n'$ is increasing. If $y'-x'=k^2$ with $k\ge1$, then
$y-x=Rk^2=(\sqrt R\,k)^2$, which is excluded.

(c) The set $\Xi_{H,D}$ depends only on the prime divisors of $D$. For
$\xi\in\Xi_{H,D}$, $R$ is coprime to $q(\xi)$, so multiplication by $R$
permutes the reduced fractions of denominator $q(\xi)$. Since $\sqrt R$ is
coprime to $q(\xi)$, the map $m\mapsto\sqrt Rm$ permutes the units modulo
$q(\xi)$, and $g(R\xi)=\varphi(q(\xi))^{-1}\sum_me(\xi(\sqrt Rm)^2)=g(\xi)$.

(d) By (a), $\widehat F_N(\xi)=N^{-1}\sum_{n'}F'(n')e(-(a_R+Rn')\xi)
=(N'/N)e(-a_R\xi)\widehat{F'}_{N'}(R\xi)$, and likewise for $J$. In the product
$\widehat F_N(-\xi)\widehat J_N(\xi)$ the phases $e(a_R\xi)e(-a_R\xi)$ cancel.
Substitute $\xi'=R\xi$, use (c), and write $D=RD_R$.
\end{proof}

The factor $R(N'/N)^2$ in (d) is essential: omitting it would
multiply the right-hand side by $R^{-1}(N/N')^2$, a factor between
$R/4$ and $R$ when $N\ge R$, since $N/R\le N'\le2N/R$
(\cref{sec:kernel}).

\subsection{The range of the method}

The closing remark records which exponent ranges the seven
conditions (E1)--(E4), (K0)--(K2) allow, and why nonunit square
classes need no new arc estimates. The unused ranges play no role
later.

\begin{remark}[The range of the parameters]\label{rem:record:K20-range}
(1) With $\kappa=\beta$ the seven conditions of \cref{def:kernel-conditions}
read $5\beta/3<\eta/2$, $13\beta/6<\eta$, $\eta+11\beta/3<1$,
$2\eta+19\beta/6<1/2$, $3\eta<1$, $3\beta/2<\eta$ and $2\eta+3\beta/2<1$. The
first and the fourth give the window $(10\beta/3,\,1/4-19\beta/12)$ of
\cref{sec:kernel}, nonempty exactly when $\beta<3/59$, and for $\eta$ in it the
other five hold as well. For squarefree $d$ and $H=N^{\beta/2}$ this is the
setting of \cite[Prop.~6.1]{OpenAI}, where $\beta=10^{-5}$ and the major arcs
have exponent $1/10$. Neither this window at values of $\beta$
other than $1/20$ nor the OpenAI parameter values just named
are used in this paper.

(2) The root weight enters the proof only through its period, its supremum
and the total size of its Fourier coefficients (each equal to $D'$), and
through the local rule of \cref{lem:record:prime-power-lift}(a). This is why
nonunit square classes, including the higher $2$-adic depths of the external
index, require no new arc estimates: the root-class lemma supplies the same
local rule and keeps $D'\le D$. The nonvanishing condition is essential to
that construction; zero square classes are not covered.
\end{remark}

\paragraph{Consequences and what to read next.}
The proof assumes \cref{def:strict-square-class,def:nonvanishing-square-class,def:kernel-notation,def:kernel-conditions},
Dirichlet's approximation
theorem and the elementary facts listed above; Weyl's inequality and the
Gauss sums at odd prime powers are proved here. At the triple of
\cref{cor:kernel20} the only numerical input is the positivity of the seven
kernel slacks, item~(S6) of \cref{cert:record:analytic-stage}.

\endgroup
